% Generated draft: requires mathematical and source review.
\documentclass[11pt]{article}
\usepackage[margin=25mm]{geometry}
\usepackage{fontspec}
\setmainfont{Latin Modern Roman}
\usepackage{amsmath,amssymb,mathtools,unicode-math}
\setmathfont{Latin Modern Math}
\usepackage{enumitem,needspace,xurl,hyperref}
\setlength{\parindent}{0pt}
\setlength{\parskip}{5pt}
\setlength{\emergencystretch}{3em}
\newcommand{\N}{\mathbb{N}}
\newcommand{\Z}{\mathbb{Z}}
\newcommand{\Q}{\mathbb{Q}}
\newcommand{\R}{\mathbb{R}}
\newcommand{\C}{\mathbb{C}}
\newcommand{\F}{\mathbb{F}}
\newcommand{\A}{\mathbb{A}}
\newcommand{\PP}{\mathbb{P}}
\newcommand{\E}{\mathbb{E}}
\newcommand{\eps}{\varepsilon}
\newcommand{\dd}{\,\mathrm{d}}
\begin{document}
\section*{UnsolvedMath ID: 11100007}
\textbf{Source problem code:} catalogue entry
\subsection{Definitions and mathematical statement}
A finite loop space consists of a based space \(BX\), a space \(X\) of finite CW homotopy type, and a specified homotopy equivalence \(X\simeq\Omega BX\). The problem is to classify all such loop-space structures up to the intended equivalence of deloopings/loop spaces, not merely classify their underlying spaces. The source does not restrict to connected \(X\).

Hovey describes a long-term program of Bill Dwyer and Clarence Wilkerson, involving Weyl groups and maximal tori, and states an expectation that Lie groups are essentially the only examples. This is an expectation, not a hypothesis permitting all other cases to be discarded. It expressly notes prime-specific exotic examples, naming \(BD_3\) at \(p=2\); that notation is not defined here.

For comparison, a connected \(p\)-compact group has a \(p\)-complete classifying space \(BY\), \(Y\simeq\Omega BY\), and finite-dimensional \(H^*(Y;\F_p)\). Classification after prime completion is not by itself a classification of all integral finite loop spaces: connectedness, compatible global data and equivalence conventions must also be addressed. The supplied summary labels a \(p\)-completed/split classification solved, gives conflicting or incomplete attribution, and ends mid-sentence. Its named \(DI(4)\) is not silently identified with the printed \(BD_3\), and its local classification is not promoted to a proved global Lie-only classification.
\subsection{Short English statement}
Classify all finite loop spaces, preserving their delooping structure and possible disconnected/global phenomena, while distinguishing prime-completed classification and exotic examples from the integral problem.
\subsection{Sources}
\begin{itemize}
\item Mark Hovey, \emph{Algebraic Topology Problem List}, ``Major problems,'' item 7, surviving mirror \url{https://www-users.cse.umn.edu/~tlawson/hovey/}. Recorded access 2026-07-29; extraction html-ordered-list. The source credits Bill Dwyer and Clarence Wilkerson's program and names \(BD_3\) without a separate reference.
\item The supplied summary names Kitchloo, Andersen--Grodal--M\o ller--Viruel, Aguad\'e--Broto--Kitchloo--Oliver (``Spaces with polynomial mod-p cohomology''), Kitchloo (``The classification of p-compact groups''), Dwyer--Miller--Wilkerson, Notbohm and Dwyer--Wilkerson's \(DI(4)\). These are incomplete/unverified supplied attributions, without exact publication details; the summary truncates after a proposed combination of local classifications.
\item Local UnsolvedMath registry, \texttt{data/unsolvedmath/problems.json}, permanent ID \(11100007\), exact source code \texttt{catalogue entry}; solved/SOLVED-IN-LITERATURE label is qualified in the summary as \(p\)-completed/split and connected.
\end{itemize}
\end{document}
